"鱼类资源量增加"给出了恢复结果，"鱼类增长的驱动机制"和"基础资源的承载压力变化"还要从模型内部追查。若将四大家鱼合为单一消费者功能群，并把水草、浮游植物、浮游动物和底栖饵料汇成单一资源总量，模型虽仍能匹配2025年资源量恢复至1.8倍，但无法辨认哪一层级资源率先承压，也难以为"水下荒漠"等底层生态退化现象提供解释。因此，正式模型中必须分别表示四类基础资源及其与四大家鱼的基本食性对应关系。
